% =====================================================================
%  Compact tool subsection (TOSEM) using the five real screenshots.
%  Preamble needs:
%     \usepackage{graphicx}
%     \usepackage{subcaption}   % for subfigure grid
%     \usepackage{url}          % or hyperref
%  Put the five PNGs in a ./figures/ directory next to the main .tex:
%     figures/tool_setup.png  tool_prompt.png  tool_baselines.png
%     figures/tool_ape.png    tool_results.png
% =====================================================================

\subsection{An Interactive Tool for Prompt Optimization}
\label{sec:tool}

To make our procedure reusable beyond the experiments reported above, we
implemented an interactive tool that operationalizes the full pipeline---the
prompting baselines, the backtracking-enhanced APE loop of
Algorithm~\ref{alg:approach}, and per-class evaluation---behind a graphical
interface. The tool requires no coding, infers the label set directly from
the data, and is decoupled from any specific model provider: the same
workflow runs against a locally hosted open-weight model (e.g., via Ollama)
or any OpenAI-compatible endpoint. It thus serves three purposes:
\emph{replicating} our results, \emph{transferring} the method to any
single-label text-classification task, and \emph{exposing} prompt
optimization as an observable, step-by-step process.

The interface follows the natural order of an experiment in five steps.
In \textbf{Data \& Config} (Fig.~\ref{fig:tool-setup}) the user uploads a CSV
whose first column is the text and last column the gold label; the tool
reports the class distribution, lets the user pick the backend and model,
and exposes the run hyper-parameters---train/test split, majority-voting
runs, the iteration budget $N_{\max}$, and the backtracking threshold
$X$---as sliders. In \textbf{Prompt} (Fig.~\ref{fig:tool-prompt}) the prompt
is split into a \emph{fixed part} that frames the task and fixes the output
format, and an \emph{optimizable part} holding the label definitions; only
the latter is rewritten during optimization, focusing the search on the
semantically meaningful portion of the prompt. \textbf{Baselines}
(Fig.~\ref{fig:tool-baselines}) runs zero-shot, few-shot, Chain-of-Thought
(CoT), and CoT~+~few-shot individually or together, reporting per-class
F\textsubscript{1} for each.

The \textbf{APE Optimisation} step (Fig.~\ref{fig:tool-ape}) is the core of
the tool. It can advance one \emph{interactive} iteration at a time or run
\emph{automatically} to the budget. At each iteration the engine
meta-prompts the model to improve the optimizable section, evaluates the
candidate on the test split with majority voting, ranks it by
F\textsubscript{1}, and---after $X$ consecutive non-improving steps---backtracks
to the next-best prompt. The iteration log, the colored iteration-history
strip (green~=~improvement, amber~=~backtrack), the live ``best prompt'' and
its evolving optimizable section, and the current per-class metrics are all
surfaced together, making the otherwise opaque dynamics of the search
directly observable. Finally, \textbf{Results} (Fig.~\ref{fig:tool-results})
contrasts every method run in the session by macro- and per-class
F\textsubscript{1}, making the gain of the optimized prompt over the
strongest baseline immediately legible. The tool, together with the datasets
and label definitions used in this paper, is available at
\url{<REPLACE-WITH-REPOSITORY-URL>}.

% ---- Figure 1: Setup (single, full width) -----------------------------------
\begin{figure}[t]
  \centering
  \includegraphics[width=0.86\linewidth]{figures/tool_setup.png}
  \caption{\textbf{Data \& Config.} The dataset is uploaded as a CSV (text
  first, label last); the tool infers the classes and their distribution,
  selects the backend and model, and exposes the run hyper-parameters as
  sliders.}
  \label{fig:tool-setup}
\end{figure}

% ---- Figure 2: Prompt (single, full width) ----------------------------------
\begin{figure}[t]
  \centering
  \includegraphics[width=0.96\linewidth]{figures/tool_prompt.png}
  \caption{\textbf{Prompt.} The prompt is separated into a fixed task-framing
  part and an optimizable definitions part (highlighted); APE rewrites only
  the latter.}
  \label{fig:tool-prompt}
\end{figure}

% ---- Figure 3: APE (single, the key figure) ---------------------------------
\begin{figure}[t]
  \centering
  \includegraphics[width=0.96\linewidth]{figures/tool_ape.png}
  \caption{\textbf{APE Optimisation.} Interactive or automatic execution of
  Algorithm~\ref{alg:approach}. The iteration log and colored history strip
  (green~=~improvement, amber~=~backtrack), the live best prompt and its
  evolving optimizable section, and the per-class metrics are shown together,
  exposing the optimization trajectory.}
  \label{fig:tool-ape}
\end{figure}

% ---- Figure 4: Baselines + Results side by side -----------------------------
\begin{figure}[t]
  \centering
  \begin{subfigure}[t]{0.49\linewidth}
    \centering
    \includegraphics[width=\linewidth]{figures/tool_baselines.png}
    \caption{Baselines: per-class F\textsubscript{1} per strategy.}
    \label{fig:tool-baselines}
  \end{subfigure}\hfill
  \begin{subfigure}[t]{0.49\linewidth}
    \centering
    \includegraphics[width=\linewidth]{figures/tool_results.png}
    \caption{Results: macro/per-class F\textsubscript{1} across methods.}
    \label{fig:tool-results}
  \end{subfigure}
  \caption{\textbf{Baselines and Results.} (a)~Each prompting strategy is run
  and scored per class. (b)~The session summary contrasts the baselines with
  the APE-optimized prompt; the optimized prompt yields the highest macro
  F\textsubscript{1} and improves the weakest classes.}
  \label{fig:tool-baseres}
\end{figure}
